\section{Orientifold KLRW algebra}
In this section, we review the definition and basic properties of KLR algebras and orientifold KLR (oKLR) algebras following \cite{KL09,Rou,VV11,Prz23}, with particular emphasis on their polynomial representations. We then introduce an extension of the oKLR algebra by adding red strands, analogous to the passage from KLR algebras to KLRW algebras; see \cite{KTWWY19}.


\subsection{KLR algebras}
\label{sec:KLR}
Fix once and for all an algebraically closed ground field $\kk$.
Let $\Qui = (Q_0, Q_1, s, t)$ be a loop-free quiver, where $Q_0$ denotes the set of vertices and $Q_1$ the set of arrows.  
For any $h \in Q_1$, we write $s(h)$ (resp.\ $t(h)$) for the source (resp.\ target) of the arrow $h$.  
For $i, j \in Q_0$, let $m_{i,j}$ denote the number of arrows from $i$ to $j$.

The associated (symmetric) \emph{Cartan matrix} $C = (c_{i,j})_{i,j \in Q_0}$ is defined by
\[
c_{i,i} := 2, 
\qquad 
c_{i,j} := - m_{i,j} - m_{j,i} \quad \text{for } i \neq j .
\]
The matrix $C$ determines a symmetric Kac--Moody algebra $\mathfrak{g}$ with Cartan subalgebra $\mathfrak{h}$. 

We fix a choice of root datum for $\mathfrak{g}$.  
This determines a \emph{weight lattice} $P$, which is a finitely generated free abelian group equipped with a symmetric bilinear form
\[
P \times P \to \Q, 
\qquad 
(\lambda, \mu) \mapsto \lambda \cdot \mu ,
\]
and containing \emph{simple roots} $(\alpha_i)_{i \in Q_0}$ and \emph{fundamental weights} $(\varpi_i)_{i \in Q_0}$ satisfying
\[
\alpha_i \cdot \alpha_j = c_{i,j},
\qquad
\alpha_i \cdot \varpi_j = \delta_{i,j},
\quad \text{for all } i,j \in Q_0 .
\]
A weight $\lambda\in P$ is said to be \emph{dominant} if $\lambda\in \displaystyle\bigoplus_{i\in Q_0}\N \varpi_i$. Let $\dom$ denote the set of all dominant weights. 

%Let $\{\alpha_i^\vee \mid i \in Q_0\}$) be the set of simple coroots associated with $\mathfrak{g}$. Let $\{\Lambda_i\mid i\in Q_0\}\subset \mathfrak{h}$ be the fundamental coweights satisfying \[
%  \langle  \Lambda_i,\alpha_j  \rangle = \delta_{ij}, \qquad i,j \in Q_0,
%\]
%where $\langle-,-\rangle$ is the natural pairing between $\mathfrak{h}$ and $\mathfrak{h}^*$.
%Let $P^\vee\subset \mathfrak{h}$ be the \emph{coweight lattice} generated by $\alpha_i^\vee$ and $\Lambda_i$.
%The involution $\tau$ on $Q_0$ naturally extends to $P^\vee$ by setting $\tau(\Lambda_i) = \Lambda_{\tau i}$ and $\tau(\alpha_i^\vee)=\alpha_{\tau i}^\vee$ for any $i\in Q_0$. 

The \emph{root lattice} is
\[
X := \bigoplus_{i \in Q_0} \Z \alpha_i \subset P ,
\]
and we set
\[
X^+ := \bigoplus_{i \in Q_0} \N \alpha_i \subset X .
\]
For $\alpha = \sum_{i \in Q_0} c_i \alpha_i \in X^+$, the \emph{height} of $\alpha$ is defined by
\(
\height(\alpha) := \sum_{i \in Q_0} c_i .
\)

Finally, let $\Word$ denote the set of all words in the alphabet $Q_0$.  
For $\alpha \in X^+$ of height $n$, let $\Word_\alpha \subset \Word$ consist of all words
\(
\bi = i_1 \cdots i_n
\)
such that
\[
\alpha_{\bi}:=\alpha_{i_1} + \cdots + \alpha_{i_n} = \alpha .
\]
The symmetric group $S_n$ acts on $\Word_\alpha$ by permuting the letters in the natural way. We denote by $s_i$ for $i=1,2,\ldots,n-1$ the simple reflections of $S_n$.


It is common and convenient to describe KLR algebras diagrammatically. Throughout, we assume that the underlying quiver is \emph{simply laced}, that is,
\[
m_{i,j} + m_{j,i} \leq 1 \qquad \text{for all } i \neq j .
\]
Let $\cR$ be the graded $\kk$-linear strict monoidal category generated by objects
$i$ for $i \in Q_0$, together with homogeneous morphisms
\[x=
\begin{tikzpicture}
    [anchorbase]
    \draw (0,0) \botlabel{i} -- (0,.4)\toplabel{i};
     \bdot{0,.2};
\end{tikzpicture} \colon i \to i \quad \text{of degree } 2,
\qquad \sigma=
\begin{tikzpicture}
    [anchorbase]
    \draw (0,0) \botlabel{i} \braidup (.4,.4) \toplabel{i};
    \draw (.4,0) \botlabel{j} \braidup (0,.4) \toplabel{j};
\end{tikzpicture}  \colon ij \to ji \quad \text{of degree } - \alpha_i \cdot \alpha_j,
\]
for all $i,j \in Q_0$ subject to the following generating relations:
\begin{gather}
\label{dotcross}
    \begin{tikzpicture}
        [baseline=.6]
         \draw (0,0) \botlabel{i} \braidup (.4,.4);
    \draw (.4,0) \botlabel{j} \braidup (0,.4);
    \bdot{.04,.3};
    \end{tikzpicture}
    =
     \begin{tikzpicture}
        [baseline=.6]
         \draw (0,0) \botlabel{i} \braidup (.4,.4);
    \draw (.4,0) \botlabel{j} \braidup (0,.4);
    \bdot{.36,.1};
    \end{tikzpicture}
    +\delta_{i,j}
    \begin{tikzpicture}
         [baseline=.6]
         \draw (0,0) \botlabel{i} \braidup (0,.4);
    \draw (.4,0) \botlabel{j} \braidup (.4,.4);
    \end{tikzpicture},
    \qquad
        \begin{tikzpicture}
        [baseline=.6]
         \draw (0,0) \botlabel{i} \braidup (.4,.4);
    \draw (.4,0) \botlabel{j} \braidup (0,.4);
    \bdot{.04,.1};
    \end{tikzpicture}
    =
     \begin{tikzpicture}
        [baseline=.6]
         \draw (0,0) \botlabel{i} \braidup (.4,.4);
    \draw (.4,0) \botlabel{j} \braidup (0,.4);
    \bdot{.36,.3};
    \end{tikzpicture}
    +\delta_{i,j}
    \begin{tikzpicture}
         [baseline=.2]
         \draw (0,0) \botlabel{i} \braidup (0,.4);
    \draw (.4,0) \botlabel{j} \braidup (.4,.4);
    \end{tikzpicture}, \\
    \label{doublecross}
    \begin{tikzpicture}
        [baseline=.8em]
         \draw (0,0) \botlabel{i} \braidup (.4,.4) \braidup (0,.8);
    \draw (.4,0) \botlabel{j} \braidup (0,.4) \braidup (.4,.8);
    \end{tikzpicture}
    =\delta_{i\neq j}
    \begin{tikzpicture}
        [baseline=1em]
         \draw (0,0) \botlabel{i} \braidup (0,.8);
    \draw (.4,0) \botlabel{j} \braidup (.4,.8);
    \coupon{.2,.4}{f_{ij}(y_1,y_2)};
    \end{tikzpicture}, \\
    \label{braidA}
    \begin{tikzpicture}
        [baseline=1em]
        \draw (0,0) \botlabel{i} \braidup (1,1);
    \draw (1,0) \botlabel{k} \braidup (0,1);
    \draw (.5,0) \botlabel{j} \braidup (0,.5) \braidup (.5,1);
    \end{tikzpicture}
    =
        \begin{tikzpicture}
        [baseline=1em]
        \draw (0,0) \botlabel{i} \braidup (1,1);
    \draw (1,0) \botlabel{k} \braidup (0,1);
    \draw (.5,0) \botlabel{j} \braidup (1,.5) \braidup (.5,1);
    \end{tikzpicture}
    +\delta_{i=k\neq j}
    \begin{tikzpicture}
        [baseline=1em]
        \draw (0,0) \botlabel{i} \braidup (0,1);
    \draw (.8,0) \botlabel{k} \braidup (.8,1);
    \draw (.4,0) \botlabel{j} \braidup (.4,.5) \braidup (.4,1);
    \coupon{.4,.5}{\frac{f_{ij}(y_3,y_2)-f_{ij}(y_1,y_2)}{y_3-y_1}};
    \end{tikzpicture},
\end{gather}
where we let $y_k$ denote the dot on the $k$th black strand read from
left to right and
\[X_{ij}(u,v)= \begin{cases}
  1 & i \not \leftarrow j\\
 u-v & i\leftarrow j\\
\end{cases} \qquad f_{ij}(u,v)=X_{ij}(u,v)X_{ji}(v,u)=
\begin{cases}
  1 & i \not \leftrightarrow j\\
 u-v & i\leftarrow j\\
v-u & i\to j
\end{cases}
\]
We refer to $\cR$ as the quiver Hecke category. Then we define the {\em KLR algebra} associated with $\alpha\in X^+$ to be 
\[
\cR(\alpha):=\bigoplus_{\bi,\bj\in \Word_{\alpha}} \Hom_{\cR}(\bi,\bj)
\]
 so that multiplication in $\cR(\alpha)$ corresponds to vertical composition of morphisms in $\cR$. There is an obvious isomorphism of categories $\bT:\cR \cong \cR^\op$ given by reflection with respect to the horizontal axis.



\subsection{oKLR algebras}
\label{sec:oKLR}


We refer to \cite[Definition~5.1]{VV11} and \cite[Definition~2.4]{Prz23} for the algebraic definition of the oKLR algebra.  
For our purposes, we adopt a diagrammatic description of the oKLR algebra, in the spirit of Section~\ref{sec:KLR}.  
In particular, the construction depends on the choice of a dominant weight.  
For any dominant weight $\mu\in \dom$ of the form
\(
\mu = \sum_{i \in Q_0} \mu_i \varpi_i ,
\)
we define
\begin{gather}
    \label{fi}
    f^\mu_i(u)=(-u)^{\mu_i} u^{\mu_{\tau i}}, \text{ for all }i\in Q_0.
\end{gather}


Our interpretation of the oKLR algebra comes from a module category over the quiver Hecke category $\cR$. Let $(\cC,\otimes,\one )$ be a $\kk$-linear strict monoidal category, and let $\cM$ be a $\kk$-linear category.  Let $\cEnd(\cM)$ denote the strict monoidal category of $\kk$-linear endofunctors of $\cM$.  A \emph{right action} of $\cC$ on $\cM$ is a monoidal functor $\bA \colon \cC \to \cEnd(\cM)^\rev$, where $\cD^\rev$ denotes the \emph{reverse} of a monoidal category $\cD$, where we reverse the tensor product.  Given such a right action, we also say that $\cM$ is a \emph{right module category} over $\cC$, or \emph{$\cC$-module} for short; see \cite[Section 7.1]{ENGO15}.  For objects $C \in \cC$ and $M \in \cM$, we set
\begin{equation} \label{beaver}
    M \otimes C := \bA(C)(M).
\end{equation}
We say that the $\cC$-module $\cM$ is \emph{strict} if $\bA$ is a strict monoidal functor.  Equivalently, $\cM$ is strict if
\[
    M \otimes (C_1 \otimes C_2) = (M \otimes C_1) \otimes C_2,\qquad
    \text{for all } M \in \cM,\ C_1,C_2 \in \cC,
\]
and similarly for morphisms. The notion of a \emph{presentation} of $\mathcal{C}$-modules by generators and relations is introduced in \cite[Sect.~2.2]{SSS25}, to which we refer for details. 

Let $\cM$, $\cN$ be strict $\cC$-modules.  A \emph{morphism of $\cC$-modules} from $\cM$ to $\cN$ is a pair $(\bH,\omega)$, where $\bH$ is a $\kk$-linear functor from $\cM$ to $\cN$ and $\omega$ is a natural isomorphism with components
\[
    \omega_{M,C} \colon \bH(M \otimes C) \xrightarrow{\cong} \bH(M) \otimes C,\qquad M \in \cM,\ C \in \cC,
\]
such that the diagram
\begin{equation} \label{buffalo}
    \begin{tikzcd}
        & \bH (M \otimes C \otimes D) \arrow[dl,"\omega_{M \otimes C, D}"'] \ar[dr,"\omega_{M, C \otimes D}"] &
        \\
        \bH(M \otimes C) \otimes D \arrow[rr,"\omega_{M,C} \otimes 1_D"] & & \bH(M) \otimes C \otimes D
    \end{tikzcd}
\end{equation}
commutes for all $M \in \cM$ and $C,D \in \mathcal{C}$.  (See \cite[Def.~7.2.1]{ENGO15} for a more general definition, where the module categories are not required to be strict.)  We say that a morphism of $\cC$-modules is \emph{strict} if $\omega_{M,C}$ is the identity morphism for all $M \in \cM$ and $C \in \cC$.  An \emph{equivalence of $\cC$-modules} is a morphism of $\cC$-modules that is also an equivalence of categories.

\begin{definition}
\label{defcRB}
    The iquiver Hecke category associated with a dominant weight $\mu$, denoted as $\cRB$, is the $\cR$-module generated by the morphisms
    \begin{gather*}
        \begin{tikzpicture}
            [anchorbase]
            \mirror{-.5}{0}{1};
            \draw (0,0) \botlabel{i} to[out=up,in=right] (-.5,.5) to[out=right,in=down] (0,1)\toplabel{\tau i};
        \end{tikzpicture} \colon i \to \tau i,\quad \text{ for all } i\in Q_0,
    \end{gather*}
    subject to the relations (for all $i,j\in Q_0$)
    \begin{gather}
    \label{quadratic}
        \begin{tikzpicture}
            [anchorbase]
            \mirror{-.25}{0}{1};
             \draw (0,0) \botlabel{i} to[out=up,in=right] (-.25,.25) to[out=right,in=down] (0,.5) to [out=up,in=right] (-.25,.75) to[out=right,in=down] (0,1)\toplabel{i};
        \end{tikzpicture}
        =
         \begin{tikzpicture}
            [anchorbase]
            \mirror{-.25}{0}{1};
             \draw (.6,0) \botlabel{i} -- (.6,1)\toplabel{i};
             \coupon{.6,.5}{f^\mu_{i}(-y_1)};
        \end{tikzpicture},
        \qquad
             \begin{tikzpicture}
            [anchorbase]
            \mirror{-.5}{0}{1};
            \draw (0,0) \botlabel{i} to[out=up,in=right] (-.5,.5) to[out=right,in=down] (0,1)\toplabel{\tau i};
            \bdot{-.1,.3};
        \end{tikzpicture}
        =-\ 
         \begin{tikzpicture}
            [anchorbase]
            \mirror{-.5}{0}{1};
            \draw (0,0) \botlabel{i} to[out=up,in=right] (-.5,.5) to[out=right,in=down] (0,1)\toplabel{\tau i};
            \bdot{-.1,.7};
        \end{tikzpicture},\\
       % -2\delta_{i,\tau i}\ 
      %   \begin{tikzpicture}
      %      [anchorbase]
      %      \mirror{-.5}{0}{1};
      %      \draw (0,0) \botlabel{i} -- (0,1)\toplabel{\tau i};
     %   \end{tikzpicture},\\
    \label{reflection}
    \begin{tikzpicture}
        [anchorbase]
        \mirror{-.25}{0}{1};
       \draw (0,0) \botlabel{i} to[out=up,in=-15] (-.25,.2) to [out=15,in=down] (.3,.6) \braidup (0,1) \toplabel{\tau i};
        \draw (.3,0) \botlabel{j} to[out=up,in=right] (-.25,.5)  to [out=right,in=down] (.3,1) \toplabel{\tau j};
    \end{tikzpicture}
    -
    \begin{tikzpicture}
        [anchorbase]
        \mirror{-.25}{0}{1};
        \draw (0,0) \botlabel{i} \braidup (.3,.4) to [out=up,in=-15] (-.25,.8) to [out=15,in=down] (0,1) \toplabel{\tau i};
        \draw (.3,0) \botlabel{j} to[out=up,in=right] (-.25,.5)  to [out=right,in=down] (.3,1) \toplabel{\tau j};
    \end{tikzpicture}
    =\ \delta_{\tau i,j}\ 
        \begin{tikzpicture}
            [anchorbase]
            \mirror{-.6}{0}{1.2};
            \draw (.4,0)\botlabel{i} \braidup (.8,.4) -- (.8,1.2) \toplabel{i};
            \draw (.8,0)\botlabel{j} \braidup (.4,.4) -- (.4,1.2) \toplabel{j};
            \coupon{.6,.7}{\frac{f^\mu_i(y_1)-f^\mu_j(y_2)}{y_1+y_2}};
        \end{tikzpicture}. 
      %   \begin{tikzpicture}
       %     [anchorbase]
      %      \mirror{-1.3}{0}{1.2};
      %      \draw (-.9,0)\botlabel{i} to[out=up,in=right] (-1.3,.2) to[out=right,in=down] (-.9,.4) -%- (-.9,1.2) \toplabel{\tau i};
     %       \draw (-.6,0)\botlabel{j}  -- (-.6,1.2)\toplabel{j};
    %        \coupon{.6,.7}{\frac{f_{ij}(y_2,-y_1)-f_{ij}(-y_2,-y_1)}{y_2}};
  %      \end{tikzpicture}, & \text{ if }i\neq \tau i,j=\tau j; \\
       %  \begin{tikzpicture}
        %    [anchorbase]
       %     \mirror{-1.3}{0}{1.2};
       %     \draw (-.9,0)\botlabel{i} to[out=up,in=right] (-1.3,.2) to[out=right,in=down] (-.9,.4) -- (-.9,1.2) \toplabel{i};
       %     \draw (-.6,0)\botlabel{j}  -- (-.6,1.2)\toplabel{j};
       %     \coupon{.4,.7}{\frac{f_{ij}(y_2,-y_1)-f_{ij}(y_2,y_1)}{y_2}};
      %  \end{tikzpicture}
      %  +
      %  \begin{tikzpicture}
      %      [anchorbase]
%\mirror{-1.3}{0}{1.2};
     %       \draw (-.7,0)\botlabel{i}  -- (-.7,1.2) \toplabel{i};
    %        \draw (-.4,0)\botlabel{j}  -- (-.4,1.2)\toplabel{j};
     %       \coupon{.4,.7}{\frac{f_{ij}(y_2,-y_1)-f_{ij}(y_2,y_1)}{y_1y_2}};
    %    \end{tikzpicture}
    %    , & \text{ if }i=\tau i,j=\tau j,i\neq j; \\
%        0 & \text{ else }.
    \end{gather}
We refer to $\lambda$ as the framing weight of $\cRB$.
\end{definition} 
\begin{remark}
    Note that our convention for \cref{reflection} differs from that in \cite{VV11,Prz23} by a sign.  
We adopt this choice so that the polynomial representation of $\cRB$ is compatible with that of the KLRW algebras; see \cref{Polklrw1,Polklrw2}.  
This sign discrepancy ultimately comes from the fact that the polynomial representation used in \cite[Proposition~2.7]{Prz23} differs by a sign from that in \cite[Lemma~4.12]{Web17a} on the generators
\(
\begin{tikzpicture}[anchorbase]
    \draw (0,0) \botlabel{i} \braidup (.4,.4) \toplabel{i};
    \draw (.4,0) \botlabel{i} \braidup (0,.4) \toplabel{i};
\end{tikzpicture}
\)
for all $i\in Q_0$.
\end{remark}

\begin{remark}
\label{embed}
In $\cRB$, the identity morphism of the unit object $\one$ is represented by the diagram
\[
\begin{tikzpicture}[anchorbase]
    \mirror{-.5}{0}{1};
\end{tikzpicture}
\; : \; \one \to \one ,
\]
which we refer to as a {mirror}.
 The relations \cref{quadratic,reflection} are diagrammatic interpretations of relations in \cite[Def. 5.1]{VV11}. %Moreover, it follows from \cref{basis} that $\cR$ can be identified with a natural subcategory of $\cRB$ whose morphism spaces consist of string diagrams in which no strand touches the mirror.
\end{remark}
For any $\alpha\in X^+$, we define $\talpha:=\tau \alpha+\alpha$ and  
\begin{gather}
\label{objRB}
    \objRB{\alpha}:=\{ \bi\in \Word\mid \talpha=\talpha_{\bi} \}.
\end{gather}
Then the oKLR algebra $\cRB(\alpha)$ associated with $\alpha$ and $\lambda$ is defined to be 
\begin{gather}
\label{oKLRalgebra}
\cRB(\alpha):= \bigoplus_{\bi,\bj\in \objRB{\alpha}}\Hom_{\cRB}(\bi,\bj).
\end{gather}
Note that $\Word_{\alpha}\subset \objRB{\alpha}$. By \cref{embed}, we may regard $\cR(\alpha)$ as a natural subalgebra of $\cRB(\alpha)$. 
\begin{lemma}
    There is a  $\kk$-linear isomorphism of categories \[\tilde{\bT}:\cRB \rightarrow \cRB^{\op}\] given by reflection with respect to the horizontal axis.
\end{lemma}
\begin{proof}
  Since reflection across the horizontal axis induces a $\kk$-linear isomorphism of categories
\(
\cR \cong \cR^{\op},
\)
it suffices to verify that the relations in \cref{quadratic,reflection} are preserved under $\tilde{\bT}$.  
This is straightforward for \cref{quadratic}.  
For \cref{reflection}, assume that $\tau i = j$. Then necessarily $i \neq j$, and the claim follows from \cref{dotcross}.
\end{proof}

%\begin{definition} {\rm (\cite[Def. 5.1]{VV11},\cite[Def. 2.4]{Prz23})}
%    Let $\beta\in \fixroot$ with $n=\height(\beta)/2$ and $\lambda=\sum_{i\in Q_0}\lambda_i \varpi_i\in P$ be a framing weight, then the oKLR  algebra $\oKLR{\beta}$ associated with $\lambda$ is the unital algebra generated by $\unit{\nu}$, $\nu\in \fixroot$, $x_k$, $1\leq k \leq n$ and $\sigma_l$, $0\leq 0\leq n-1$, subjecting to the following relations:
%\end{definition}


%\[
% \begin{tikzpicture}
%\mirror{-.25}{0}{1};
%\mirror{4.25}{1}{0};
% \draw (0,0) to[out=up,in=right] (-.5,.5) to[out=right,in=down] (0,1);
%\draw (0,0) to[out=up,in=-15] (-.25,.5);  
%\draw (-.25,.5) to[out=15,in=down] (0,1);
%\end{tikzpicture}
%\]

%\[\begin{tikzpicture}[yscale=0.75,baseline=2em]
%\mirror{-0.25}{0}{2};
%\draw(0,0) \braidup (0.5,1) to[out = up, in = -15] (-.25,1.75) to[out = 15, in=down] (0,2);
%\draw(0.5,0) to[out = up, in = -15] (-.25,1) to[out = 15, in=down] (0.5,2);
%\end{tikzpicture}
%=begin{tikzpicture}[yscale=0.75,baseline=2em]
%\mirror{-0.25}{0}{2};
%\draw(0,0)  to[out = up, in = -15] (-.25,0.25)  to[out = 15, in=down] (0.5,1) \braidup (0,2);
%\draw(0.5,0) to[out = up, in = -15] (-.25,1) to[out = 15, in=down] (0.5,2);
% \draw(0.5,0) \braidup (0,0.65) to[out = up, in = -15] (-.25,1) to[out = 15, in=down] (0,1.35)\braidup (0.5,2);
%\end{tikzpicture}\]


\subsection{KLRW algebras}
We now recall the KLRW algebra following \cite[Definition~4.1]{Web17a}.  
In \cite{Web17a}, it is defined as the $\kk$-span of Stendhal diagrams modulo certain local relations. For our purposes, we instead present the definition in the language of strict monoidal categories, in parallel with the approach used for $\cR$.
\begin{definition} {\rm (\cite[Def.~4.1,Def.~4.4]{Web17a})}
    \label{KLRWdef}
    The KLRW category $\klrw$ is the graded $\kk$-linear strict monoidal category generated by objects $i$ and $\rlam$ for all $i\in Q_0$ and all $\rlam\in \dom$ together with homogeneous morphisms
    \begin{gather*}
\begin{tikzpicture}
    [anchorbase]
    \draw (0,0) \botlabel{i} -- (0,.4)\toplabel{i};
     \bdot{0,.2};
\end{tikzpicture} \colon i \to i \quad \text{of degree } 2,
\qquad 
\begin{tikzpicture}
    [anchorbase]
    \draw (0,0) \botlabel{i} \braidup (.4,.4) \toplabel{i};
    \draw (.4,0) \botlabel{j} \braidup (0,.4) \toplabel{j};
\end{tikzpicture}  \colon ij \to ji \quad \text{of degree } - \alpha_i \cdot \alpha_j, \\
\begin{tikzpicture}
    [anchorbase]
    \draw (0,0) \botlabel{i} \braidup (.4,.4) \toplabel{i};
    \draw[red,thick] (.4,0) \botlabel{\rlam} \braidup (0,.4) \toplabel{\rlam};
\end{tikzpicture}  \colon i\rlam \to \rlam i \quad \text{of degree }  \alpha_i \cdot \rlam, \quad
\begin{tikzpicture}
    [anchorbase]
    \draw[red,thick] (0,0) \botlabel{\rlam} \braidup (.4,.4) \toplabel{\rlam};
    \draw (.4,0) \botlabel{i} \braidup (0,.4) \toplabel{i};
\end{tikzpicture}  \colon \rlam i \to i\rlam  \quad \text{of degree }  \alpha_i \cdot \rlam,
    \end{gather*}
for all $i,j \in Q_0$ and all dominant weights $\rlam$, subjecting to \cref{dotcross}--\cref{braidA}  and the following generating relations:
\begin{gather}
\label{redbraid1}
\begin{tikzpicture}
        [baseline=1em]
        \draw (0,0) \botlabel{i} \braidup (1,1);
    \draw (1,0) \botlabel{j} \braidup (0,1);
    \draw[red,thick] (.5,0) \botlabel{\rlam} \braidup (0,.5) \braidup (.5,1);
    \end{tikzpicture}
    ~=~
        \begin{tikzpicture}
        [baseline=1em]
        \draw (0,0) \botlabel{i} \braidup (1,1);
    \draw (1,0) \botlabel{j} \braidup (0,1);
    \draw[red,thick] (.5,0) \botlabel{\rlam} \braidup (1,.5) \braidup (.5,1);
    \end{tikzpicture}
    ~+\delta_{i,j}\sum_{a+b+1=\rlam_i}~
     \begin{tikzpicture}
        [baseline=1em]
        \draw (0,0) \botlabel{i} \braidup (0,1);
        \bdot{0,.5};
        \bdot{1,.5};
        \node at (-.2,.5){$\scriptstyle a$};
        \node at (1.2,.5){$\scriptstyle b$};
    \draw (1,0) \botlabel{j} \braidup (1,1);
    \draw[red,thick] (.5,0) \botlabel{\rlam} \braidup (.5,.5) \braidup (.5,1);
    \end{tikzpicture}, \\
\label{redbraid2}
\begin{tikzpicture}
        [baseline=1em]
        \draw (0,0) \botlabel{i} \braidup (1,1);
    \draw[red,thick] (1,0) \botlabel{\rlam} \braidup (0,1);
    \draw (.5,0) \botlabel{j} \braidup (0,.5) \braidup (.5,1);
    \end{tikzpicture}
    ~=~
        \begin{tikzpicture}
        [baseline=1em]
        \draw (0,0) \botlabel{i} \braidup (1,1);
    \draw[red,thick] (1,0) \botlabel{\rlam} \braidup (0,1);
    \draw (.5,0) \botlabel{j} \braidup (1,.5) \braidup (.5,1);
    \end{tikzpicture}, \quad
    \begin{tikzpicture}
        [baseline=1em]
        \draw[red,thick] (0,0) \botlabel{\rlam} \braidup (1,1);
    \draw (1,0) \botlabel{i} \braidup (0,1);
    \draw (.5,0) \botlabel{j} \braidup (0,.5) \braidup (.5,1);
    \end{tikzpicture}
    ~=~
        \begin{tikzpicture}
        [baseline=1em]
        \draw[red,thick] (0,0) \botlabel{\rlam} \braidup (1,1);
    \draw (1,0) \botlabel{i} \braidup (0,1);
    \draw (.5,0) \botlabel{j} \braidup (1,.5) \braidup (.5,1);
    \end{tikzpicture}, \\
    \label{dotmovered}
    \begin{tikzpicture}
        [baseline=.6]
         \draw (0,0)[red,thick] \botlabel{\rlam} \braidup (.4,.4);
    \draw (.4,0) \botlabel{i} \braidup (0,.4);
    \bdot{.04,.3};
    \end{tikzpicture}
    =
     \begin{tikzpicture}
        [baseline=.6]
         \draw (0,0)[red,thick] \botlabel{\rlam} \braidup (.4,.4);
    \draw (.4,0) \botlabel{i} \braidup (0,.4);
    \bdot{.36,.1};
    \end{tikzpicture},\quad 
        \begin{tikzpicture}
        [baseline=.6]
         \draw (0,0) \botlabel{i} \braidup (.4,.4);
    \draw (.4,0)[red,thick] \botlabel{\rlam} \braidup (0,.4);
    \bdot{.04,.1};
    \end{tikzpicture}
    =
     \begin{tikzpicture}
        [baseline=.6]
         \draw (0,0)\botlabel{i} \braidup (.4,.4);
    \draw (.4,0)[red,thick]  \botlabel{\rlam} \braidup (0,.4);
    \bdot{.36,.3};
    \end{tikzpicture}, \\
    \label{reddoublecross}
    \begin{tikzpicture}
        [baseline=.8em]
         \draw (0,0)[red,thick] \botlabel{\rlam} \braidup (.4,.4) \braidup (0,.8);
    \draw (.4,0) \botlabel{i} \braidup (0,.4) \braidup (.4,.8);
    \end{tikzpicture}
    = \begin{tikzpicture}
        [baseline=.8em]
         \draw (0,0)[red,thick] \botlabel{\rlam} \braidup (0,.4) \braidup (0,.8);
    \draw (.4,0) \botlabel{i} \braidup (.4,.4) \braidup (.4,.8);
    \bdot{.4,.4}; \node at (.6,.4){$\scriptstyle \rlam_i$};
    \end{tikzpicture}, \quad
    \begin{tikzpicture}
        [baseline=.8em]
         \draw (0,0) \botlabel{i} \braidup (.4,.4) \braidup (0,.8);
    \draw (.4,0)[red,thick] \botlabel{\rlam} \braidup (0,.4) \braidup (.4,.8);
    \end{tikzpicture}
    =
     \begin{tikzpicture}
        [baseline=.8em]
         \draw (0,0) \botlabel{i} \braidup (0,.4) \braidup (0,.8);
         \bdot{0,.4}; \node at (-.2,.4){$\scriptstyle \rlam_i$};
    \draw (.4,0)[red,thick] \botlabel{\rlam} \braidup (.4,.4) \braidup (.4,.8);
    \end{tikzpicture},
\end{gather}
where $\rlam=\sum_{i\in Q_0}\rlam_i \varpi_i$ and $i,j\in Q_0$.
\end{definition}

Let $\obj{\klrw}$ denote the set of objects of $\klrw$.  
In particular, each object in $\obj{\klrw}$ is uniquely determined by a triple
\(
(\bi,\redkappa,\lamseq),
\)
where $\bi=(i_1,\dots,i_n)$ records the labels of the black strands from left to right, $\redkappa$ is a weakly increasing function
\(
\redkappa \colon [1,\ell] \to [0,n]
\)
such that the $k$-th red strand lies between the $\redkappa(k)$-th and $(\redkappa(k)+1)$-st black strands, and
\[
\lamseq = (\rlam^{(1)}, \rlam^{(2)}, \ldots, \rlam^{(\ell)})
\]
labels the red strands from left to right.  
We interpret $\redkappa(k)=0$ to mean that the $k$-th red strand lies to the left of all black strands, and $\redkappa(k)=n$ to mean that it lies to the right of all black strands.  
When $\lamseq$ is clear from the context, we omit it and write simply $(\bi,\redkappa)$ for the corresponding object. For any two sequences of dominant weights $\lamseq$ and $\lamseq\red{'}$, we write $\lamseq \cup \lamseq\red{'}$ for their obvious horizontal concatenation. Moreover, for any $\lamseq$ we define 
\begin{gather}
\label{objklrwalg}
\obj{\klrw}^{\lamseq}=\{ (\bi,\redkappa,\red{\underline{\lambda'}})\in \obj{\klrw}\mid \red{\underline{\lambda'}}=\lamseq \}.
\end{gather}
Then we define the KLRW algebra associated with $\lamseq$ to be 
\begin{equation}
    \label{klrwalg}
    \klrw^{\lamseq}:=\oplus_{(\bi,\redkappa_1),(\bj,\redkappa_2)\in \obj{\klrw}^{\lamseq}}\Hom_{\klrw}((\bi,\redkappa_1),(\bj,\redkappa_2)).
\end{equation}
We note that in \cref{klrwalg} $\Hom_{\klrw}((\bi,\redkappa_1),(\bj,\redkappa_2))=\{0\}$ unless $\alpha_\bi=\alpha_\bj$.

\begin{remark}
    As Webster explained in \cite[Prop. 4.21]{Web17a}, one can split $\lamseq$ into a sequence consisting of fundamental weights only, denoted by $\lamseq^{\red{\operatorname{spl}}}$, then the algebra $\klrw^{\lamseq}$ is isomorphic to an idempotent truncation of $\klrw^{\lamseq^{\red{\operatorname{spl}}}}$.
\end{remark}

As in \cite[Def. 4,7]{Web17a}, we let $\klrw_{\alpha}^\lamseq$ be the subalgebra of $\klrw^{\lamseq}$ where the sum of the roots associated to the black strands is $\sum_{i=1}^\ell\rlam_i-\alpha$ and let $\klrw_{n}^\lamseq$ denote the subalgebra of diagrams with $n$ black strands.

We then describe the polynomial representation of $\klrw^{\lamseq}$ inherited from the one of $\klrw$. Define
\begin{equation}
\label{eq:Pol}
 \Pol(\bi,\redkappa,\lamseq):=\kk[Y_1(\bi,\redkappa,\lamseq),\dots Y_n(\bi,\redkappa,\lamseq)] \text{ where }\bi=\{i_1,\ldots,i_n\}.
\end{equation}
As implied by \cite[Lem. 4.12]{Web17a}, there is a strict monoidal functor $\cP: \klrw \to \cVect_\kk$ such that
\begin{equation}
\label{Polobj}
\cP(\bi,\kappa,\lamseq)= \Pol(\bi,\kappa,\lamseq)
\end{equation}
for all $(\bi,\kappa,\lamseq)\in \obj{\klrw}$, and 
\begin{gather}
\label{Polklrw1}
    \cP\Big(\begin{tikzpicture}
        [baseline=.6]
         \draw (0,0)[red,thick] \botlabel{\rlam} \braidup (.4,.4);
    \draw (.4,0) \botlabel{i} \braidup (0,.4);
    \end{tikzpicture}\Big)
    ~: f\mapsto Y_1^{\rlam_i}\cdot f~\quad
    \cP\Big(\begin{tikzpicture}
        [baseline=.6]
         \draw (0,0) \botlabel{i} \braidup (.4,.4);
    \draw (.4,0)[red,thick] \botlabel{\rlam} \braidup (0,.4);
    \end{tikzpicture}\Big)
    ~: f\mapsto f~\quad
    \cP\Big(\begin{tikzpicture}
        [baseline=.6]
        \draw (0,0)\botlabel{i} -- (0,.4);
        \bdot{0,.2};
    \end{tikzpicture}\Big)
    ~: f\mapsto Y_1\cdot f \\
    \label{Polklrw2}
    \cP\Big(\begin{tikzpicture}
        [baseline=.3]
         \draw (0,0) \botlabel{i} \braidup (.6,.6);
    \draw (.6,0) \botlabel{j} \braidup (0,.6);
    \end{tikzpicture}\Big)
    ~: f\mapsto \begin{cases}
        X_{ij}(Y_{2},Y_1)f^{s_1}, & i\neq j, \\[1ex]
        \dfrac{f^{s_1}-f}{Y_{2}-Y_{1}}, & i=j.
    \end{cases}
\end{gather}
where 
\begin{align}
\label{weylact}
f^{s_j}(Y_1,\ldots,Y_n) =
\begin{cases} 
 f(\cdots,Y_{j+1},Y_{j},\cdots),&\text{ if } j>0; \\
 f(-Y_1,Y_2,\cdots), &\text{ if } j=0.
\end{cases} 
\end{align}
In particular, restricting $\cP$ to $\klrw^{\lamseq}$ defines a polynomial representation of $\klrw^{\lamseq}$ on \[\Pol^\lamseq:=\oplus_{(\bi,\redkappa)\in \obj{\klrw}^{\lamseq}}\Pol(\bi,\redkappa).\]

\subsection{oKLRW algebras}
\label{sec:oKLRW}
We are now ready to introduce the oKLRW category and the oKLRW algebra. Recall the iquiver Hecke category $\cRB$ from \cref{defcRB} as a module category over $\cR$. 
\begin{definition}
    \label{defoKLRWC}
    The oKLRW category associated with a dominant weight $\mu$, denoted as $\oklrwmu$, is the $\klrw$-module generated by the morphisms
     \begin{gather*}
        \begin{tikzpicture}
            [anchorbase]
            \mirror{-.5}{0}{1};
            \draw (0,0) \botlabel{i} to[out=up,in=right] (-.5,.5) to[out=right,in=down] (0,1)\toplabel{\tau i};
        \end{tikzpicture} \colon i \to \tau i,\quad \text{ for all } i\in Q_0,
    \end{gather*}
    subject to the relations \cref{quadratic,reflection} (which depend on $\mu$).
\end{definition}
By \cref{defoKLRWC}, we observe that the red strands will never touch the mirror. Moreover, the category $\oklrwmu$ shares the same set of objects as $\klrw$. Let $\lamseq$ be a fixed sequence of dominant weights, recall $\obj{\klrw}^{\lamseq}$ from \cref{objklrwalg}. We define the oKLRW algebra to be
\begin{gather}
\label{eq:oklrwalg}
    \oklrwmu^{\lamseq}:=\bigoplus_{(\bi,\redkappa_1),(\bj,\redkappa_2)\in \obj{\klrw}^{\lamseq}} \Hom_{\oklrwmu}((\bi,\redkappa_1),(\bj,\redkappa_2)).    
\end{gather}
As in \cref{embed}, we can regard the KLRW algebra $\klrwalg$ as a natural subalgebra of the oKLRW algebra $\oklrwalg$. Recall the functor $\cP$ from \cref{Polobj,Polklrw1,Polklrw2}.
\begin{prop}
\label{prop:polyrep}
     Regarding $\cVect_\kk$ as a strict $\klrw$-module via $\cP$,  there is a unique strict morphism of $\klrw$-modules $\ctP:\oklrwmu\to \cVect_\kk$ such that
     \begin{gather}
         \label{Polreflect}
         \ctP \Big(~
            \begin{tikzpicture}
            [anchorbase]
            \mirror{-.5}{0}{1};
            \draw (0,0) \botlabel{i} to[out=up,in=right] (-.5,.5) to[out=right,in=down] (0,1)\toplabel{\tau i};
        \end{tikzpicture} 
        ~ \Big) : f \mapsto (-Y_1)^{\mu_i}f^{s_0}.
     \end{gather}
\end{prop}
\begin{proof}
    By our definition, $\ctP$ agrees with $\cP$ when restricting to $\klrw$. Thus it suffices to show that \cref{Polreflect} satisfies \cref{quadratic,reflection}, which is straightforward to check; see also \cite[Prop 2.7]{Prz23}.
\end{proof}
To summarize, we now have the following diagram:
\begin{equation}
\label{diagram}
\begin{tikzcd}
&\text{Quiver Hecke category } \cR \arrow[rd] \arrow[ld]& \\
\text{iQuiver Hecke category } \cRB \arrow[rd,dashed] & & \text{KLRW category } \klrw  \arrow[ld] \\
& \text{oKLRW category } \oklrwmu&
\end{tikzcd}
\end{equation}
To interpret the dashed arrow in \cref{diagram}, we set $\mu=\bzero$ and consider the full subcategory
\[
{}_{\bzero}\cS_\lambda \subset \oklrwzero
\]
whose objects are those containing exactly one red strand labelled by
\(
\lamseq = (\lambda).
\)
 By definition, ${}_\bzero\cS_\lambda$ is naturally a $\cR$-module.
\begin{theorem}
\label{framezero}
    For any framing weight $\lambda$, 
    there exists an embedding of $\cR$-modules  $\cF\colon\cRB\hookrightarrow {}_\bzero\cS_\lambda$ which sends any object $\bi\in \cRB$ to $(\bi,\redkappa,\rlam)$ with $\redkappa(1)=0$ and
    \begin{gather}
        \begin{tikzpicture}
            [anchorbase]
            \mirror{-.5}{0}{1};
            \draw (0,0) \botlabel{i} to[out=up,in=right] (-.5,.5) to[out=right,in=down] (0,1)\toplabel{\tau i};
        \end{tikzpicture} ~\mapsto~~\begin{tikzpicture}
            [anchorbase]
            \mirror{-.5}{0}{1};
            \draw (0,0) \botlabel{i} to[out=up,in=right] (-.5,.5) to[out=right,in=down] (0,1)\toplabel{\tau i};
            \draw[red] (-.25,0)\botlabel{\mu} \braidup (.1,.5) \braidup (-.25,1);
        \end{tikzpicture} 
    \end{gather}
\end{theorem}
%The following proposition shows that the iquiver Hecke category $\cRB$ can be regarded as a subcategory of $\oklrwzero$.
\begin{proof}
    We observe that $\ctP\circ\cF$ defines a faithful polynomial representation of $\cRB$ which agrees with \cite[Prop. 2.7]{Prz23} (up to a sign), hence $\cF$ is faithful as well. Then it suffices to check that \cref{quadratic,reflection} are preserved under $\cF$. For the first relation in \cref{quadratic}, in ${}_\bzero\cS_\lambda$ we have
    \begin{gather*}
        \begin{tikzpicture}
            [anchorbase]
            \mirror{-.5}{0}{1};
             \draw (0,0) \botlabel{i} to[out=up,in=right] (-.5,.25) to[out=right,in=down] (0,.5) to [out=up,in=right] (-.5,.75) to[out=right,in=down] (0,1)\toplabel{i};
             \draw[red] (-.25,0) \botlabel{\mu} \braidup (0,.25) \braidup (-.25,.5) \braidup (0,.75)\braidup (-.25,1);
        \end{tikzpicture}
        \overset{\cref{reddoublecross}}{=}
         \begin{tikzpicture}
            [anchorbase]
            \mirror{-.5}{0}{1};
             \draw (0,0) \botlabel{i} to[out=up,in=right] (-.5,.25) to[out=right,in=down] (-.25,.5) to [out=up,in=right] (-.5,.75) to[out=right,in=down] (0,1)\toplabel{i};
             \draw[red] (-.25,0) \botlabel{\mu} \braidup (.2,.5) \braidup (-.25,1);
              \bdot{-.25,.5}; \node at (0,.5){$\scriptstyle \rmu_{\tau i}$};
        \end{tikzpicture}
        \underset{\cref{dotmovered}}{\overset{\cref{quadratic}}{=}}(-1)^{\rmu_{\tau i}}
        \begin{tikzpicture}
            [anchorbase]
            \mirror{-.5}{0}{1};
             \draw (0,0) \botlabel{i} to[out=up,in=right] (-.5,.25) to[out=right,in=down] (-.25,.5) to [out=up,in=right] (-.5,.75) to[out=right,in=down] (0,1)\toplabel{i};
             \draw[red] (-.25,0) \botlabel{\mu} \braidup (0,.5) \braidup (-.25,1);
              \bdot{-.02,.9}; \node at (.2,.8){$\scriptstyle \rmu_{\tau i}$};
        \end{tikzpicture}
        \overset{\cref{quadratic}}{=}(-1)^{\rmu_{\tau i}}
        \begin{tikzpicture}
            [anchorbase]
            \mirror{-.5}{0}{1};
             \draw (0,0) \botlabel{i} \braidup (-.25,.5)\braidup (0,1)\toplabel{i};
             \draw[red] (-.25,0) \botlabel{\mu} \braidup (0,.5) \braidup (-.25,1);
              \bdot{-.02,.9}; \node at (.2,.8){$\scriptstyle \rmu_{\tau i}$};
        \end{tikzpicture}
        \overset{\cref{reddoublecross}}{=}(-1)^{\rmu_{\tau i}}
         \begin{tikzpicture}
            [anchorbase]
            \mirror{-.5}{0}{1};
             \draw (0,0) \botlabel{i} \braidup  (0,1)\toplabel{i};
             \draw[red] (-.25,0) \botlabel{\mu} \braidup (-.25,1);
              \bdot{0,.7}; \node at (.3,.7){$\scriptstyle \rmu_{\tau i}$};
               \bdot{0,.3}; \node at (.3,.3){$\scriptstyle \rmu_{i}$};
        \end{tikzpicture}
        =   \begin{tikzpicture}
            [anchorbase]
            \mirror{-.25}{0}{1};
            \draw[red] (-.1,0)\botlabel{\mu} -- (-.1,1);
             \draw (.6,0) \botlabel{i} -- (.6,1)\toplabel{i};
             \coupon{.6,.5}{f^{\rmu}_{i}(-y_1)};
        \end{tikzpicture}.
    \end{gather*}
Thus, $\cF$ preserves the first relation in \cref{quadratic}.
For the second relation in \cref{quadratic}, \cref{dotmovered} implies that it is preserved by $\cF$ as well.

For \cref{reflection}, first we suppose that $i\neq \tau j$, and hence $\tau i\neq j$. Then we have 
\begin{gather*}
        \begin{tikzpicture}
        [anchorbase]
        \mirror{-.25}{0}{1};
       \draw (0,0) \botlabel{i} to[out=up,in=-15] (-.25,.2) to [out=15,in=down] (.3,.6) \braidup (0,1) \toplabel{\tau i};
        \draw (.3,0) \botlabel{j} to[out=up,in=right] (-.25,.5)  to [out=right,in=down] (.3,1) \toplabel{\tau j};
        \draw[red] (-.125,0)\botlabel{\mu} \braidup (0,.3) \braidup (-.125,1);
    \end{tikzpicture}
    \overset{\cref{redbraid1}}{=}
     \begin{tikzpicture}
        [anchorbase]
        \mirror{-.25}{0}{1};
       \draw (0,0) \botlabel{i} to[out=up,in=-15] (-.25,.2) to [out=15,in=down] (.3,.6) \braidup (0,1) \toplabel{\tau i};
        \draw (.3,0) \botlabel{j} to[out=up,in=right] (-.25,.5)  to [out=right,in=down] (.3,1) \toplabel{\tau j};
        \draw[red] (-.125,0)\botlabel{\mu} \braidup (.5,.3) \braidup (-.125,1);
    \end{tikzpicture}
    \overset{\cref{redbraid2}}{=}
     \begin{tikzpicture}
        [anchorbase]
        \mirror{-.25}{0}{1};
       \draw (0,0) \botlabel{i} to[out=up,in=-15] (-.25,.2) to [out=15,in=down] (.3,.6) \braidup (0,.8) \braidup (.2,1) \toplabel{\tau i};
        \draw (.3,0) \botlabel{j} to[out=up,in=right] (-.25,.5)  to [out=right,in=down] (.5,1) \toplabel{\tau j};
        \draw[red] (-.125,0)\botlabel{\mu} \braidup (.5,.5) \braidup (0,1);
    \end{tikzpicture}
    \overset{\cref{reflection}}{=}
     \begin{tikzpicture}
        [anchorbase]
        \mirror{-.25}{0}{1};
      \draw (0,0) \botlabel{i} \braidup (-.125,.2) \braidup (.3,.4) to [out=up,in=-15] (-.25,.8) to [out=15,in=down] (.2,1) \toplabel{\tau i};
        \draw (.3,0) \botlabel{j} to[out=up,in=right] (-.25,.5)  to [out=right,in=down] (.5,1) \toplabel{\tau j};
        \draw[red] (-.125,0)\botlabel{\mu} \braidup (.5,.5) \braidup (0,1);
    \end{tikzpicture}
    \underset{\cref{redbraid2}}{\overset{\cref{redbraid1}}{=}}
    \begin{tikzpicture}
        [anchorbase]
        \mirror{-.25}{0}{1};
        \draw (0,0) \botlabel{i} \braidup (.3,.4) to [out=up,in=-15] (-.25,.8) to [out=15,in=down] (0,1) \toplabel{\tau i};
        \draw (.3,0) \botlabel{j} to[out=up,in=right] (-.25,.5)  to [out=right,in=down] (.5,1) \toplabel{\tau j};
         \draw[red] (-.125,0)\botlabel{\mu} \braidup (0,.3) \braidup (-.125,1);
    \end{tikzpicture}.
\end{gather*}
Now suppose that $i=\tau j$, hence $i\neq j$. Then in ${}_\bzero\cS_\lambda\subset \oklrwzero$ we have
\begin{gather}
\label{steak1}
      \begin{tikzpicture}
        [anchorbase]
        \mirror{-.25}{0}{1};
       \draw (0,0) \botlabel{i} to[out=up,in=-15] (-.25,.2) to [out=15,in=down] (.3,.6) \braidup (0,1) \toplabel{j};
        \draw (.3,0) \botlabel{j} to[out=up,in=right] (-.25,.5)  to [out=right,in=down] (.3,1) \toplabel{i};
        \draw[red] (-.125,0)\botlabel{\mu} \braidup (0,.3) \braidup (-.125,1);
    \end{tikzpicture}
    \underset{\cref{redbraid2}}{\overset{\cref{redbraid1}}{=}}
    \begin{tikzpicture}
        [anchorbase]
        \mirror{-.25}{0}{1};
       \draw (0,0) \botlabel{i} to[out=up,in=-15] (-.25,.2) to [out=15,in=down] (.3,.6) \braidup (0,.8) \braidup (.2,1) \toplabel{j};
        \draw (.3,0) \botlabel{j} to[out=up,in=right] (-.25,.5)  to [out=right,in=down] (.5,1) \toplabel{i};
        \draw[red] (-.125,0)\botlabel{\mu} \braidup (.5,.5) \braidup (0,1);
    \end{tikzpicture}
    ~+\sum_{a+b+1=\rmu_j}~
    \begin{tikzpicture}
        [anchorbase]
        \mirror{-.5}{0}{1};
        \draw (-.1,0) \botlabel{i} to[out=up,in=-15] (-.5,.3) to [out=15,in=down] (-.3,.5) to[out=up,in=-15] (-.5,.7) to [out=15,in=down] (.1,1) \toplabel{i};
        \draw (.1,0) \botlabel{j} \braidup (-.1,1) \toplabel{j};
        \draw[red] (-.3,0)\botlabel{\mu} \braidup (-.1,.5) \braidup (-.3,1);
        \bdot{-.3,.5}; \node at (-.18,.5){$\scriptstyle a$};
        \bdot{.08,.2}; \node at (.2,.2){$\scriptstyle b$};
    \end{tikzpicture}
    \underset{\cref{dotmovered}}{\overset{\cref{quadratic},\cref{reddoublecross}}{=}}
     \begin{tikzpicture}
        [anchorbase]
        \mirror{-.25}{0}{1};
       \draw (0,0) \botlabel{i} to[out=up,in=-15] (-.25,.2) to [out=15,in=down] (.3,.6) \braidup (0,.8) \braidup (.2,1) \toplabel{j};
        \draw (.3,0) \botlabel{j} to[out=up,in=right] (-.25,.5)  to [out=right,in=down] (.5,1) \toplabel{i};
        \draw[red] (-.125,0)\botlabel{\mu} \braidup (.5,.5) \braidup (0,1);
    \end{tikzpicture}
    ~+\sum_{a+b+1=\rmu_j}(-1)^a~
    \begin{tikzpicture}
        [anchorbase]
        \mirror{-1}{0}{1};
        \draw (.1,0) \botlabel{i} \braidup (.5,1) \toplabel{i};
        \draw (.5,0) \botlabel{j} \braidup (.1,1) \toplabel{j};
        \draw[red] (-.8,0)\botlabel{\mu} \braidup (-.8,1);
        \bdot{.48,.2}; \node at (.6,.2){$\scriptstyle b$};
        \bdot{.12,.2}; \node at (-.3,.38){$\scriptstyle \rmu_i+a$};
    \end{tikzpicture}
\\
\label{steak2}
    \begin{tikzpicture}
        [anchorbase]
        \mirror{-.25}{0}{1};
        \draw (0,0) \botlabel{i} \braidup (.3,.4) to [out=up,in=-15] (-.25,.8) to [out=15,in=down] (0,1) \toplabel{j};
        \draw (.3,0) \botlabel{j} to[out=up,in=right] (-.25,.5)  to [out=right,in=down] (.5,1) \toplabel{i};
         \draw[red] (-.125,0)\botlabel{\mu} \braidup (0,.3) \braidup (-.125,1);
    \end{tikzpicture}
    \underset{\cref{redbraid2}}{\overset{\cref{redbraid1}}{=}}
     \begin{tikzpicture}
        [anchorbase]
        \mirror{-.25}{0}{1};
      \draw (0,0) \botlabel{i} \braidup (-.125,.2) \braidup (.3,.4) to [out=up,in=-15] (-.25,.8) to [out=15,in=down] (.2,1) \toplabel{j};
        \draw (.3,0) \botlabel{j} to[out=up,in=right] (-.25,.5)  to [out=right,in=down] (.5,1) \toplabel{i};
        \draw[red] (-.125,0)\botlabel{\mu} \braidup (.5,.5) \braidup (0,1);
    \end{tikzpicture}
    ~+\sum_{c+d+1=\rmu_i}~
     \begin{tikzpicture}
        [anchorbase]
        \mirror{-.5}{0}{1};
        \draw (.1,0) \botlabel{j} to[out=up,in=-15] (-.5,.3) to [out=15,in=down] (-.3,.5) to[out=up,in=-15] (-.5,.7) to [out=15,in=down] (-.1,1) \toplabel{j};
        \draw (-.1,0) \botlabel{i} \braidup (.1,1) \toplabel{i};
        \draw[red] (-.3,0)\botlabel{\mu} \braidup (-.1,.5) \braidup (-.3,1);
        \bdot{-.3,.5}; \node at (-.18,.5){$\scriptstyle c$};
        \bdot{0,.5}; \node at (.2,.5){$\scriptstyle d$};
    \end{tikzpicture}
    \underset{\cref{dotmovered},\cref{reddoublecross}}{\overset{\cref{dotcross},\cref{quadratic}}{=}}
    \begin{tikzpicture}
        [anchorbase]
        \mirror{-.25}{0}{1};
      \draw (0,0) \botlabel{i} \braidup (-.125,.2) \braidup (.3,.4) to [out=up,in=-15] (-.25,.8) to [out=15,in=down] (.2,1) \toplabel{j};
        \draw (.3,0) \botlabel{j} to[out=up,in=right] (-.25,.5)  to [out=right,in=down] (.5,1) \toplabel{i};
        \draw[red] (-.125,0)\botlabel{\mu} \braidup (.5,.5) \braidup (0,1);
    \end{tikzpicture}
    ~+\sum_{c+d+1=\rmu_i}(-1)^c~
    \begin{tikzpicture}
        [anchorbase]
        \mirror{-.3}{0}{1};
        \draw (.1,0) \botlabel{i} \braidup (.5,1) \toplabel{i};
        \draw (.5,0) \botlabel{j} \braidup (.1,1) \toplabel{j};
        \draw[red] (-.2,0)\botlabel{\mu} \braidup (-.2,1);
        \bdot{.48,.2}; \node at (.9,.2){$\scriptstyle \rmu_j+c$};
        \bdot{.12,.2}; \node at (-.05,.3){$\scriptstyle d$};
    \end{tikzpicture}
\end{gather}
Notice that
\begin{gather} 
    \label{steak3}
    \sum_{a+b+1=\rmu_j}(-1)^a~
    \begin{tikzpicture}
        [anchorbase]
        \mirror{-1}{0}{1};
        \draw (.1,0) \botlabel{i} \braidup (.5,1) \toplabel{i};
        \draw (.5,0) \botlabel{j} \braidup (.1,1) \toplabel{j};
        \draw[red] (-.8,0)\botlabel{\mu} \braidup (-.8,1);
        \bdot{.48,.2}; \node at (.6,.2){$\scriptstyle b$};
        \bdot{.12,.2}; \node at (-.3,.38){$\scriptstyle \rmu_i+a$};
    \end{tikzpicture}
    = 
    \begin{tikzpicture}
            [anchorbase]
            \mirror{-.8}{0}{1.2};
            \draw[red] (-.7,0)\botlabel{\mu} -- (-.7,1.2);
            \draw (.4,0)\botlabel{i} \braidup (.4,.8) \braidup (.8,1.2) \toplabel{i};
            \draw (.8,0)\botlabel{j} \braidup (.8,.8) \braidup (.4,1.2) \toplabel{j};
            \coupon{.6,.5}{\frac{y_2^{\rmu_j}-(-y_1)^{\rmu_j}}{y_1+y_2}y_1^{\rmu_i}};
        \end{tikzpicture},\quad
    \sum_{c+d+1=\rmu_i}(-1)^c~
    \begin{tikzpicture}
        [anchorbase]
        \mirror{-.3}{0}{1};
        \draw (.1,0) \botlabel{i} \braidup (.5,1) \toplabel{i};
        \draw (.5,0) \botlabel{j} \braidup (.1,1) \toplabel{j};
        \draw[red] (-.2,0)\botlabel{\mu} \braidup (-.2,1);
        \bdot{.48,.2}; \node at (.9,.2){$\scriptstyle \rmu_j+c$};
        \bdot{.12,.2}; \node at (-.05,.3){$\scriptstyle d$};
    \end{tikzpicture}
    =
     \begin{tikzpicture}
            [anchorbase]
            \mirror{-.8}{0}{1.2};
            \draw[red] (-.7,0)\botlabel{\mu} -- (-.7,1.2);
            \draw (.4,0)\botlabel{i} \braidup (.4,.8) \braidup (.8,1.2) \toplabel{i};
            \draw (.8,0)\botlabel{j} \braidup (.8,.8) \braidup (.4,1.2) \toplabel{j};
            \coupon{.6,.5}{\frac{y_1^{\rmu_i}-(-y_2)^{\rmu_i}}{y_1+y_2}y_2^{\rmu_j}};
        \end{tikzpicture}
\end{gather}
Since  $ \begin{tikzpicture}
        [anchorbase]
        \mirror{-.25}{0}{1};
       \draw (0,0) \botlabel{i} to[out=up,in=-15] (-.25,.2) to [out=15,in=down] (.3,.6) \braidup (0,.8) \braidup (.2,1) \toplabel{j};
        \draw (.3,0) \botlabel{j} to[out=up,in=right] (-.25,.5)  to [out=right,in=down] (.5,1) \toplabel{i};
        \draw[red] (-.125,0)\botlabel{\mu} \braidup (.5,.5) \braidup (0,1);
    \end{tikzpicture}
    =
    \begin{tikzpicture}
        [anchorbase]
        \mirror{-.25}{0}{1};
      \draw (0,0) \botlabel{i} \braidup (-.125,.2) \braidup (.3,.4) to [out=up,in=-15] (-.25,.8) to [out=15,in=down] (.2,1) \toplabel{j};
        \draw (.3,0) \botlabel{j} to[out=up,in=right] (-.25,.5)  to [out=right,in=down] (.5,1) \toplabel{i};
        \draw[red] (-.125,0)\botlabel{\mu} \braidup (.5,.5) \braidup (0,1);
    \end{tikzpicture}$ in ${}_\bzero\cS_\lambda$,
by taking difference of \cref{steak1,steak2}, we obtain that
\begin{gather*}
     \begin{tikzpicture}
        [anchorbase]
        \mirror{-.25}{0}{1};
       \draw (0,0) \botlabel{i} to[out=up,in=-15] (-.25,.2) to [out=15,in=down] (.3,.6) \braidup (0,1) \toplabel{j};
        \draw (.3,0) \botlabel{j} to[out=up,in=right] (-.25,.5)  to [out=right,in=down] (.3,1) \toplabel{i};
        \draw[red] (-.125,0)\botlabel{\mu} \braidup (0,.3) \braidup (-.125,1);
    \end{tikzpicture}
    -
     \begin{tikzpicture}
        [anchorbase]
        \mirror{-.25}{0}{1};
        \draw (0,0) \botlabel{i} \braidup (.3,.4) to [out=up,in=-15] (-.25,.8) to [out=15,in=down] (0,1) \toplabel{j};
        \draw (.3,0) \botlabel{j} to[out=up,in=right] (-.25,.5)  to [out=right,in=down] (.5,1) \toplabel{i};
         \draw[red] (-.125,0)\botlabel{\mu} \braidup (0,.3) \braidup (-.125,1);
    \end{tikzpicture}
    \overset{\cref{steak3}}{=}
    \begin{tikzpicture}
            [anchorbase]
            \mirror{-.8}{0}{1.2};
            \draw[red] (-.65,0)\botlabel{\mu} -- (-.65,1.2);
            \draw (.4,0)\botlabel{i} \braidup (.4,.8) \braidup (.8,1.2) \toplabel{i};
            \draw (.8,0)\botlabel{j} \braidup (.8,.8) \braidup (.4,1.2) \toplabel{j};
            \coupon{.6,.5}{\frac{f^{\rmu}_i(y_2)-f^{\rmu}_j(y_1)}{y_1+y_2}};
        \end{tikzpicture}
        \overset{\cref{dotcross}}{=}
        \begin{tikzpicture}
            [anchorbase]
            \mirror{-.8}{0}{1.2};
            \draw[red] (-.65,0)\botlabel{\mu} -- (-.65,1.2);
            \draw (.4,0)\botlabel{i} \braidup (.8,.4) -- (.8,1.2) \toplabel{i};
            \draw (.8,0)\botlabel{j} \braidup (.4,.4) -- (.4,1.2) \toplabel{j};
            \coupon{.6,.7}{\frac{f^{\rmu}_i(y_1)-f^{\rmu}_j(y_2)}{y_1+y_2}};
        \end{tikzpicture}
\end{gather*}
This finishes the proof.
\end{proof}
\cref{framezero} shows that, for any framing weight $\lambda$, we can always regard $\cRB$ as a subcategory of $\oklrwzero$ by adding a red strand labeled $\lambda$ to the left end of each generating morphism. The exact same proof shows the following:
\begin{corollary}
\label{embedzeroklrw}
    For any framing weight $\mu$, there exists an embedding of $\klrw$-modules
\(
\tilde{\cF}\colon \oklrwmu \hookrightarrow \oklrwzero,
\)
which sends an object $(\bi,\redkappa,\lamseq)$ to
\(
(\bi,\redkappa',\, \mu \cup \lamseq),
\)
where \(\redkappa'(1)=0, 
\redkappa'(k+1)=\redkappa(k)\ \text{ for all } k\ge 1,
\)
and such that
    \begin{gather*}
        \begin{tikzpicture}
            [anchorbase]
            \mirror{-.5}{0}{1};
            \draw (0,0) \botlabel{i} to[out=up,in=right] (-.5,.5) to[out=right,in=down] (0,1)\toplabel{\tau i};
        \end{tikzpicture} ~\mapsto~~\begin{tikzpicture}
            [anchorbase]
            \mirror{-.5}{0}{1};
            \draw (0,0) \botlabel{i} to[out=up,in=right] (-.5,.5) to[out=right,in=down] (0,1)\toplabel{\tau i};
            \draw[red] (-.25,0)\botlabel{\mu} \braidup (.1,.5) \braidup (-.25,1);
        \end{tikzpicture} 
    \end{gather*}
\end{corollary}
Therefore, by \cref{framezero,embedzeroklrw}, to study the oKLRW category associated with an arbitrary framing weight, it suffices to understand the case where the framing weight is $\bzero$.

\subsection{Diagonal type}
\label{sec:diagnoal}
Recall the notion of a quiver with an involution from \cref{sec:oKLR}. In this subsection, we assume that our quiver $\Qui=(Q_0,Q_1,s,t)$ is the disjoint union of two
 copies of quivers with identical vertices and opposite arrows, denoted by
\[
\Qui=\Qui_{+}\bigsqcup \Qui_{-}
\]
and let $\tau :\Qui\to\Qui$  be the involution that interchanges
the vertices in corresponding positions. For example, we may have the following quiver with involution
$$
\begin{tikzpicture}[
  scale=.8,
  every node/.style={circle, draw, inner sep=3pt},
  arrow/.style={red, <->, thick}
]

\node (t1) at (-1,1) {};   % upper fork
\node (t2) at (0,0) {};   % center
\node (t3) at (-1,-1) {};  % lower fork
\node (t4) at (1,0) {};
\node (t5) at (2,0) {};
\node (t6) at (3,0) {};
\node (t7) at (4,0) {};

\draw[->,ultra thick] (t1) -- (t2); 
\draw[->,ultra thick] (t2) -- (t3);
\draw[<-,ultra thick] (t2) -- (t4);
\draw[<-,ultra thick] (t4) -- (t5);
\draw[<-,ultra thick] (t5) -- (t6);
\draw[->,ultra thick] (t6) -- (t7);

\node (b1) at (-1,0) {};   % lower fork
\node (b2) at (0,-1) {};   % center
\node (b3) at (-1,-2) {};   % upper fork
\node (b4) at (1,-1) {};
\node (b5) at (2,-1) {};
\node (b6) at (3,-1) {};
\node (b7) at (4,-1) {};


\draw[<-,ultra thick] (b1) -- (b2);
\draw[<-,ultra thick] (b2) -- (b3);
\draw[->,ultra thick] (b2) -- (b4);
\draw[->,ultra thick] (b4) -- (b5);
\draw[->,ultra thick] (b5) -- (b6);
\draw[<-,ultra thick] (b6) -- (b7);

\draw[arrow] (t1) to[bend right=30] (b1);
\draw[arrow] (t2) to[bend left=20] (b2);
\draw[arrow] (t3) to[bend right=30] (b3);
\draw[arrow] (t4) to[bend left=15] (b4);
\draw[arrow] (t5) to[bend left=15] (b5);
\draw[arrow] (t6) to[bend left=15] (b6);
\draw[arrow] (t7) to[bend left=15] (b7);

\end{tikzpicture}$$
where each copy is a type D Dynkin quiver.

Let $\lamseq$ be a fixed sequence of dominant weights associated with $\Qui_{+}$. Recall the oKLRW algebra ${\oklrwzero}^{\lamseq}$ associated with $(\Qui,\tau)$ with framing weight $\bzero$ from \cref{oKLRalgebra} and the KLRW algebra ${\klrwalg}$ associated with $\Qui^{(+)}$ from \cref{klrwalg}.
\begin{theorem}
\label{thm:diagonal}
    The oKLRW algebra ${\oklrwzero}^{\lamseq}$ associated with $(\Qui,\tau)$ is Morita equivalent to the KLRW algebra ${\klrwalg}$ associated with $\Qui_{+}$. 
\end{theorem}
\begin{proof}
    Let $Q_{0,\pm}$ denote the set of indices of $\Qui_{\pm}$. By our definition, $\klrwalg$ is a span of Stendhal diagrams (\cite[Def. 4.1]{Web17a}) and elements of $\obj{\klrw}^{\lamseq}$ are of the form $(\bi,\redkappa)$ where $\bi$ are words in $Q_{0,+}$. Let \[
    e:=\sum_{(\bi,\redkappa)\in \obj{\klrw}^{\lamseq}} \id_{(\bi,\redkappa)} \in {\oklrwzero}^{\lamseq}
    \]
 where $\id_{(\bi,\redkappa)}$ is the identity morphism of $(\bi,\redkappa)$ within ${\oklrwzero}$. To show that ${\oklrwzero}^{\lamseq}$ is Morita equivalent to ${\klrwalg}$, it suffices to show 
 \[
 e({\oklrwzero}^{\lamseq})e\cong \klrwalg \quad\text{ and }\quad ({\oklrwzero}^{\lamseq} )e({\oklrwzero}^{\lamseq})={\oklrwzero}^{\lamseq}.
 \]
 Note that 
 \[
   e({\oklrwzero}^{\lamseq})e=\oplus_{(\bi,\redkappa_1),(\bj,\redkappa_2)\in \obj{\klrw}^{\lamseq}} \Hom_{\oklrwzero}((\bi,\redkappa_1),(\bj,\redkappa_2)).    
 \]
Since $\bi,\bj$ are both words in $Q_{0,+}$, we see that any black strand in those Stendhal diagrams must touch the mirror even number of times. However, inside $\oklrwzero$, by \cref{quadratic,reflection} we have for any $i,j$
\[
 \begin{tikzpicture}
            [anchorbase]
            \mirror{-.25}{0}{1};
             \draw (0,0) \botlabel{i} to[out=up,in=right] (-.25,.25) to[out=right,in=down] (0,.5) to [out=up,in=right] (-.25,.75) to[out=right,in=down] (0,1)\toplabel{i};
        \end{tikzpicture}
        =
         \begin{tikzpicture}
            [anchorbase]
            \mirror{-.25}{0}{1};
             \draw (.1,0) \botlabel{i} -- (.1,1)\toplabel{i};
             %\coupon{.6,.5}{f^\lambda_{i}(-y_1)};
        \end{tikzpicture}
    \text{ and }
     \begin{tikzpicture}
        [anchorbase]
        \mirror{-.25}{0}{1};
       \draw (0,0) \botlabel{i} to[out=up,in=-15] (-.25,.2) to [out=15,in=down] (.3,.6) \braidup (0,1) \toplabel{\tau i};
        \draw (.3,0) \botlabel{j} to[out=up,in=right] (-.25,.5)  to [out=right,in=down] (.3,1) \toplabel{\tau j};
    \end{tikzpicture}
    =
    \begin{tikzpicture}
        [anchorbase]
        \mirror{-.25}{0}{1};
        \draw (0,0) \botlabel{i} \braidup (.3,.4) to [out=up,in=-15] (-.25,.8) to [out=15,in=down] (0,1) \toplabel{\tau i};
        \draw (.3,0) \botlabel{j} to[out=up,in=right] (-.25,.5)  to [out=right,in=down] (.3,1) \toplabel{\tau j};
    \end{tikzpicture}.
\]
Thus we may always use the second relation above, together with \cref{braidA,redbraid1,redbraid2}, to detach the black strands from the mirror.  
It follows that $e(\oklrwzero^{\lamseq})e$ is spanned by Stendhal diagrams in which the black strands never touch the mirror and whose domains and codomains lie in $\obj{\klrw}^{\lamseq}$.  
Consequently, the identification with $\klrwalg$ is obtained simply by removing the mirror; see also \cref{embed}.

For the second equality, it suffices to show that for every object $(\bi,\redkappa,\lamseq)$ in $\oklrwzero$, the identity morphism
\(
\id_{(\bi,\redkappa,\lamseq)}
\)
lies in the ideal $(\oklrwzero^{\lamseq})\,e\,(\oklrwzero^{\lamseq})$.  
Here $\bi$ is a word in the alphabet $Q_0$, and it may contain letters in $Q_{0,-}$.  
We argue by an induction on the number $k$ of occurrences of vertices from $Q_{0,-}$ in the word $\bi$.
When $k=0$, there is nothing to prove. Suppose the statement is true for $k-1$, and $i_l$ is the $k$-th occurrence of vertices from $Q_{0,-}$ in the word $\bi$. Note that we may cross any red strand with the black strand labelled by $i_l$ freely since weights in $\lamseq$ are associated with $\Qui_{+}$ only. By the inductive hypothesis, we may assume $i_1,i_2,\ldots i_{l-1}\in  Q_{0,+}$ and then we have (omitting the positions of the red strands):
\begin{gather*}
   \begin{tikzpicture}
       [anchorbase]
       \mirror{-.25}{0}{1};
       \draw (0,0)\botlabel{i_1} -- (0,1)\toplabel{};
       \draw (.25,0)\botlabel{i_2} -- (.25,1);
       \node at (.5,.5){$\scriptstyle \cdots$};
       \draw (.75,0)\botlabel{i_{l-1}} -- (.75,1);
       \draw (1.25,0)\botlabel{i_l} -- (1.25,1);
      \node at (1.5,.5){$\scriptstyle \cdots$};
   \end{tikzpicture}
   \overset{\cref{doublecross}}{=}
   \begin{tikzpicture}
       [anchorbase]
       \mirror{-.25}{0}{1};
       \draw (0,0)\botlabel{i_1} -- (0,1)\toplabel{};
       \draw (.25,0)\botlabel{i_2} -- (.25,1);
       \node at (.5,.5){$\scriptstyle \cdots$};
       \draw (.75,0)\botlabel{i_{l-1}} -- (.75,1);
       \draw (1.25,0)\botlabel{i_l} to[out=150,in=down] (-.15,.5) to[out=up,in=-150] (1.25,1);
      \node at (1.5,.5){$\scriptstyle \cdots$};
    %  \draw[dashed] (-.25,.65) -- (1.5,.65);
     % \draw[dashed] (-.25,.35) -- (1.5,.35);
   \end{tikzpicture}
   \overset{\cref{reflection}}{=}
      \begin{tikzpicture}
       [anchorbase]
       \mirror{-.25}{0}{1};
       \draw (0,0)\botlabel{i_1} -- (0,1)\toplabel{};
       \draw (.25,0)\botlabel{i_2} -- (.25,1);
       \node at (.5,.5){$\scriptstyle \cdots$};
       \draw (.75,0)\botlabel{i_{l-1}} -- (.75,1);
       \draw (1.25,0)\botlabel{i_l} to[out=150,in=-15] (-.25,.3) to [out=15,in=-15] (-.25,.7) to[out=15,in=-150] (1.25,1);
      \node at (1.5,.5){$\scriptstyle \cdots$};
      \draw[dashed] (-.25,.65) -- (1.5,.65);
      \draw[dashed] (-.25,.35) -- (1.5,.35);
   \end{tikzpicture}
   \in (\oklrwzero^{\lamseq})\,e\,(\oklrwzero^{\lamseq})
\end{gather*}
where the middle identity morphism is a summand of $e$ since $\tau i_l\in Q_0^{(+)}$. This completes the proof.
\end{proof}

\subsection{Basis}
Recall the Weyl group of type $B_n$ is isomorphic to the signed symmetric group 
$SS_n$, the group of permutations $w$ of 
$\{\pm 1, \pm 2, \ldots, \pm n\}$ such that 
$w(-i)=-w(i)$. It is generated by 
$\{s_i\mid 0\le n-1 \}$, where $s_0=(-1,1)$ and $s_i=(i,i+1)(-i,-i-1)$. 

We will use  the mirror strand 
diagrams to represent each $w$ in $SS_n$. These are a special kind of $2n$-strand diagrams such that each strand connects $i$ and $w(i)$, but the picture is required to be symmetric across a central axis (the mirror). For example, 
\[ s_0= 
\begin{tikzpicture}
            [anchorbase]
            \mirror{-.5}{0}{1};
            \draw (0,0) \botlabel{1} to[out=up,in=right] (-.5,.5) to[out=right,in=down] (0,1)\toplabel{1};
            \draw (-1,0) \botlabel{-1} to[out=up,in=left] (-.5,.5) to[out=left,in=down] (-1,1)\toplabel{-1};
            \draw (.5,0)\botlabel{2} to (.5,1)\toplabel{2};
            \node at (1,.5) {$\cdots$};
            \draw (1.5,0)\botlabel{n} to (1.5,1)\toplabel{n};
             \draw (-1.5,0)\botlabel{-2} to (-1.5,1)\toplabel{-2};
            \node at (-2,.5) {$\cdots$};
            \draw (-2.5,0)\botlabel{-n} to (-2.5,1)\toplabel{-n};
        \end{tikzpicture}.
        \]
  So, each mirror pair of crossings corresponds to a simple reflection $s_i$, while a crossing which occurs on the top of the mirror corresponds to the simple refection $s_0$. 

  Since the picture is symmetric, we can only draw the right half of the picture so as to corresponds to the diagrams in $\cRB$. For example, 
  \[s_0= 
  \begin{tikzpicture}
            [anchorbase]
            \mirror{-.5}{0}{1};
            \draw (0,0) \botlabel{1} to[out=up,in=right] (-.5,.5) to[out=right,in=down] (0,1)\toplabel{1};
            \draw (.5,0)\botlabel{2} to (.5,1)\toplabel{2};
            \node at (1,.5) {$\cdots$};
            \draw (1.5,0)\botlabel{n} to (1.5,1)\toplabel{n};
        \end{tikzpicture}.
        \]
 For any  $\bi\in Q_0^n$ and $w\in SS_n$, let $w\bi=(j_1,\ldots,j_n)\in Q_0^n$ be obtained from 
 $\bi$ such that if $w(k)>0$,  then $j_{w(k)}=i_k$ and if $w(k)<0$, then $j_{-w(k)}=\tau(i_k) $. 

  We fix a sequence of dominant weight $\lamseq$ in the following. 
\begin{definition}
\label{def:reduceddiagram}
For any $w\in SS_n$ and  $(\bi,\redkappa)$
and $\redkappa':[1,\ell]\rightarrow [0,n]$, we  choose a Stendhal diagram (without dots)  $D_{w,\redkappa'}=1_{w\bi, \redkappa',\lamseq}D_{w,\redkappa'} 1_{\bi, \redkappa,\lamseq}$  from $ (\bi,\redkappa,\lamseq)$ to $(w\bi, \redkappa',\lamseq)$ such that 
\begin{enumerate}
    \item the underlying diagram without red lines is  a mirror strand  diagram of $w$
    such that no two black strands cross twice and each strand meets the mirror at most once. In other words, the crossing reading bottom to top is a reduced expression for $w$.
    \item the red and black strand which does not meet the mirror cross at most once. 
    \item the red and black strands which meet the mirror cross at most twice.
    We call such a diagram a  reduced Stendhal diagram for $w$ and $(\bi,\redkappa)$ and $\redkappa'$. 
\end{enumerate}
\end{definition}

\begin{example} For the following two diagrams, 
\[ D=\begin{tikzpicture}
            [anchorbase]
            \mirror{-.5}{-1}{1};
            \draw (0.3,-1)to [out=up,in=down](.5,-.5)  to[out=up,in=right] (-.5,.5) to[out=right,in=down] (0,1);
            \draw (.6,1) to [out=down,in=right] (-.5,0) to [out=right,in=up] (.9,-1);
            \draw (.6,-1) to (1.5,1);
            \draw (0.9,1) to (1.5,-1);
            \draw[red] (.2,1) to (-.2,-1);
              \draw[red] (1.2,1) to [out=down, in=up] (1.5,0) to [out=down,in=up] (1.2,-1);
                     \end{tikzpicture}, \quad 
       D'=\begin{tikzpicture}
            [anchorbase]
            \mirror{-.5}{-1}{1};
            \draw (0,-1)to [out=up,in=down](.3,-.5)  to[out=up,in=right] (-.5,.5) to[out=right,in=down] (0,1);
            \draw (.6,1) to [out=down,in=right] (-.5,0) to [out=right,in=up] (.9,-1);
            \draw (.6,-1) to (1.5,1);
            \draw (0.9,1) to (1.5,-1);
            \draw[red] (-.2,1) to   (.2,-1);
              \draw[red] (1.2,1) to [out=down, in=up] (1,0) to [out=down,in=up] (1.7,-1);
                     \end{tikzpicture}              \]
we have  that $D$ is a reduced Stendhal diagram. While $D'$ satisfies condition (1)—as its underlying diagram (omitting red lines) corresponds to a reduced expression for some $w \in SS_n$—it fails to satisfy (2) and (3). Specifically, the first red line intersects a black strand meeting the mirror three times, while the second red line intersects a non-mirror-meeting black strand twice.   Thus, $D'$ is not a  reduced Stendhal diagram.    
\end{example}

Let 
$B_{\redkappa',\redkappa}=\{ D_{w,\redkappa'}y^\mathbf a1_{\bi, \redkappa,\lamseq} \mid w\in SS_n, \bi\in Q_0^n, \mathbf a\in \N^n \}$, 
where $y^\mathbf a=y_1^{a_1}\ldots y_n^{a_n}$.

\begin{lemma}
\label{lem:spanpart}
    The set $B=\bigcup_{\redkappa',\redkappa}B_{\redkappa',\redkappa}$ 
    spans $\oklrwmu^{\lamseq}$.
\end{lemma}
\begin{proof}
We prove by induction on the number $c$ of crossings that any diagram $D$ can be written as a linear combination of elements in $B$. Note that we also view the meet between the strand and the mirror as a crossing since it represents $s_0$ in the language of the mirror strand diagram. 

If $c=0$, then $D$ is the identity morphism of some $1_{\bi,\redkappa,\lamseq} $ with some dots and hence $D\in B$. Suppose $c>0$.
By \eqref{dotcross}, \eqref{quadratic}, and \eqref{dotmovered}, dots can be moved freely through crossings modulo diagrams with fewer crossings. So, we can assume  the dots of $D$ are all at the bottom.  
Using the quadratic \eqref{doublecross}, \eqref{reddoublecross} and \eqref{quadratic}, and also 
 the braid  relations \eqref{braidA} and \eqref{reflection}, and \eqref{redbraid1}--\eqref{redbraid2},  we can assume that  
(1)--(3) in Definition \ref{def:reduceddiagram} hold for the underlying diagram dropping the dots of $D$. Finally, the braid relations also show that for any two choices $D'_{w,\redkappa'}$ and $D_{w,\redkappa'}$ for any  fixed $w$, we must have $D '_{w,\redkappa'}=D_{w,\redkappa'} $  modulo diagrams with fewer crossings. This completes the proof by inductive hypotheses.
\end{proof}

The proof of the following result follows the same triangularity argument as the KLRW basis theorem of Webster, combined with the PBW basis for oKLR algebras.

\begin{theorem}
\label{oklrwbasis}
   The set  $B=\bigcup_{\redkappa',\redkappa}B_{\redkappa',\redkappa}$ 
    is a basis of  $\oklrwmu^{\lamseq}$.
\end{theorem}
\begin{proof}By Lemma \ref{lem:spanpart}, it remains to show that $B$ is linearly independent.
  We first prove the claim that $\{D_{w,\red{n^\ell}}y^\mathbf a 1_{\bi, \red{n^\ell},\lamseq} \mid w\in SS_n, \bi\in Q_0^n, \mathbf a\in \N^n \}$ is linearly independent.
  In fact, we consider the obvious  functor from 
$ _\mu ^\tau\mathcal R$ to  $\oklrwmu^{\lamseq}$, sending each diagram $b$ to the diagram obtained by adding the 
 $\ell$ red lines to the right of $b$.
 Then composing this functor with the polynomial representation $\ctP$ yields the faithful polynomial representation of $ _\mu ^\tau\mathcal R$ in \cite{Prz23}. Hence the claim follows from the PBW basis of  $ _\mu ^\tau\mathcal R$ in \cite[Proposition 2.9]{Prz23} since $B_{\red{n^\ell},\red{n^\ell}}$ is the image of the corresponding  basis element of $ _\mu ^\tau\mathcal R$.

 Next,  we define for any $\redkappa$ an element $\theta_\redkappa$ such that it is a sum over all $\bi$ of the unique   Stendhal  diagram from  $ (\bi,\red{n^\ell}, \lamseq  ) $ to $ (\bi,\redkappa, \lamseq ) $ which has no dots and a minimal number of crossings. 
 Let $\bar \theta_\redkappa$ be the diagram obtained from $\theta_\redkappa$ by reflecting the diagram with respect to the horizontal axis. Then by the 
 same arguments as in \cite[Lemma 4.17]{Web15}, the diagram $\bar \theta_{\redkappa'}D_{w,\redkappa'}y^{\mathbf{a}} 1_{\bi,\redkappa,\lamseq} \theta_\redkappa$ is equal to 
 $D_{w,\red{n^\ell}} y^{\mathbf{a}+\mathbf{b}} 1_{\bi, \red{n^\ell},\lamseq} $ modulo the span of diagrams with fewer crossings than $D_{w, \red{n^\ell}}$, for some  $\mathbf{b}\in \mathbb Z_{\ge0}^n$ determined by the relations in \eqref{reddoublecross}.

 Now we are ready to prove the linear independence of $B_{\redkappa',\redkappa}$. Suppose
 there is a non-trivial linear relation $\sum c_{w,\mathbf{a}}D_{w,\redkappa'}y^{\mathbf{a}} 1_{\bi,\redkappa,\lamseq}=0$ with $c_{w,\mathbf a }\neq 0$.
 Multiply $\bar \theta_{\redkappa'}$ on the left and $\theta_\redkappa$ on the right. Then using the same argument as in the proof of  Lemma \ref{lem:spanpart}, each term of the summation can be rewritten as elements of form $D_{w',\red{n^\ell}}y^{\mathbf a'} 1_{\bi,\red{n^\ell},\lamseq}$.  Furthermore, choose a $w\in SS_n$ and $\mathbf a$ in the summation  (i.e., $c_{w,\mathbf a}\neq 0$) such that $w$ is maximal in the Bruhat order. Then we multiply  by $\bar \theta_{\redkappa'}$ on the left and $\theta_\redkappa$ on the right and rewrite in  terms of elements in $B_{\red{n^\ell},\red{n^\ell}}$.  The claim in the previous paragraph shows that $D_{w,\red{n^\ell}}y^{\mathbf a+\mathbf b}1_{\bi,\red{n^\ell},\lamseq} $ also has coefficient $c_{w,\mathbf a}$. since no element of  $B_{\red{n^\ell},\red{n^\ell}}$ other than  $D_{w',\red{n^\ell}}y^{\mathbf a'} $ could contribute to its coefficient. Since $B_{\red{n^\ell},\red{n^\ell}}$ is linearly independent (which was shown in the beginning of the proof), we must have $c_{w,\mathbf a}=0$, a contradiction. This show the original linear relation must be trivial and hence $B$ is linearly independent.   
\end{proof}


\begin{corollary}
\label{cor:faithfulpolyrep}
The polynomial representation from \cref{prop:polyrep} is faithful.
\end{corollary}

\begin{proof}
It suffices to prove faithfulness on each morphism space. Fix two objects
\[
(\bi,\redkappa,\lamseq)
\qquad\text{and}\qquad
(\bj,\redkappa',\lamseq)
\]
with \(n\) black strands, and set
\(
K:=\operatorname{Frac}\kk[Y_1,\ldots,Y_n].
\)
After extending the polynomial representation to \(K\), the action of a
reduced Stendhal diagram has the triangular form
\[
\ctP(D_{w,\redkappa'})
=
q_w w+\sum_{u<w}q_{u,w}u,
\qquad q_w\in K^\times,
\]
where \(<\) is the Bruhat order on \(SS_n\). Indeed, this follows directly
from \cref{Polklrw1,Polklrw2,Polreflect}: each black crossing or mirror
reflection has a nonzero coefficient of the corresponding simple
reflection, while all other terms have smaller Bruhat order; red--black
crossings act by multiplication by nonzero polynomials.

By \cref{oklrwbasis}, every element of the fixed morphism space has a
unique expression
\[
x=\sum_{w\in SS_n}
D_{w,\redkappa'}p_w(y)
1_{\bi,\redkappa,\lamseq},
\qquad
p_w(y)\in\kk[y_1,\ldots,y_n].
\]
Assume that \(x\neq0\), and choose \(w\) maximal in the Bruhat order such
that \(p_w\neq0\). The coefficient of \(w\) in \(\ctP(x)\) is then
\(
q_w\,w\bigl(p_w(Y)\bigr),
\)
which is nonzero. Since distinct elements of \(SS_n\) induce linearly
independent automorphisms of \(K\), this term cannot cancel. Hence
\(\ctP(x)\neq0\), proving the result.
\end{proof}


\begin{remark}
\label{rem:extend}
    The results of this section admit straightforward generalizations to quivers with involution that have $\tau$-fixed vertices.  
The corresponding oKLR algebra and its polynomial representation have already been developed in \cite{Prz23}, and one can define the analogous oKLRW category/algebra in the same way as in \cref{defoKLRWC}. However, for simplicity we do not treat this case here.
\end{remark}